\scanpage{29}
$=$ si $A$ courbe elliptique sur corps alg. clos; donc l'isom can n'est \uline{pas} fonctoriel pour les hom de schémas abéliens de $=$ dimension [dans la mesure où ça pourrait avoir un sens, les deux foncteurs n'étant pas de $=$ variance en $A$\ldots]. Cependant, \deleted{\unsure{si $u$ est} \illegible{2} \unsure{A.B.} \illegible{2}}
\[
\boxed{\det T_p(u)_{\text{ét}}\,\det T_p(u^*)_{\text{ét}}=\deg u}
\]
\margin{Vaut aussi (via \unsure{identifications} \unsure{permises|possibles}) pour des Sch.A. de $=$ dimension}

c) Il faudrait vérifier que
\[
\alpha_{A^*}=\alpha_A^{-1}
\]
[ aux signes (?)\ldots]. Noter aussi la compatibilité par sommes de V.A.\ \ldots

 d) L'isom $({}_pA)_{\text{ét}}\longrightarrow({}_pA^*)_{\text{ét}}$ induit par $\alpha_A$ est bien celui déduit plus haut de la dualité en cohomologie de DeRham (à vérifier \ldots).

\mspage{28}
